%% fuer Drucker "bilby":
%% dvips -O 5mm,1mm

%% fuer Drucker "goofy":
%% dvips600 -O 0mm,10mm Magdeburg

\documentclass[portrait]{seminar}
\usepackage{german,colordvi,boldmath,semcolor,semlayer,rotate,epsf,fancybox,mathsym}
%\input{farbunilogo}
%\input{unilogo}
\usepackage{amssymb}
\def\dummy#1{}
\overlaystrue
\def\overlay#1{{}}

\slidewidth9.5in
\slideheight6.7in
%\def\slidetopmargin{\iflandscape.2in\else0.8in\fi}
\def\slidetopmargin{\iflandscape.3in\else0.8in\fi}
\def\slidebottommargin{\iflandscape1.0in\else0.4in\fi}
\def\slideleftmargin{\iflandscape.2in\else.5in\fi}
\def\sliderightmargin{\iflandscape1.0in\else.8in\fi}

\newcommand{\cH}{{\cal H}}

\slidesmag{4}
\centerslidesfalse
\slideframewidth0.5pt
%\slideframe{none}

%\portraitonly
%\landscapeonly
%\def\leftsliderotation#1{\rotate[l]{#1}}
%\sliderotation{left}
\rotateheaderstrue

%%%%%%%%%%%%%%%%%%%%%%%%%%%%%%%%%%%%%%%%%%%%%%
%% Makros
%%%%%%%%%%%%%%%%%%%%%%%%%%%%%%%%%%%%%%%%%%%%%%
\def\bra#1{\left<#1\right|}
\def\ket#1{\left|#1\right>}

\def\trace{\mathop{\rm tr}\nolimits}

\newcommand{\heading}[1]{%
  \begin{center}
    \Large
    \bfseries
    \shadowbox{#1}%
  \end{center}
  \vspace{1ex minus 1ex}%
}
\newbox{\logobox}
%\setbox\logobox=\hbox{\raise 0.17\baselineskip\hbox{\epsfbox{logo.eps}}}

\def\logo{\usebox{\logobox}}
\newpagestyle{MH}{\hss\hbox to \hsize{PhysComp96\hss\thepage}\hss}%
{\hss\hbox to \hsize{${\cal IAKS}$\ \logo\ UniKA\hfil}\hss}

\newpagestyle{nohead}{}{\hss\hbox to \hsize{\tiny IAKS
\ \logo\ Universit"at Karlsruhe\hfil Prof.~Dr.~Th.~Beth}\hss}

\newpagestyle{nohead_beth}{}{\hss\hbox to \hsize{\tiny${\cal IAKS}$\ \logo\
UniKA\hfil}\hss}

\newcommand{\tl}{\tilde}
\newcommand{\onemat}{{\mathbf 1}}

\begin{document}

\begin{slide}
{$ $}
\vfill
\begin{center}
{\Large \blue Mutual simulation of Hamiltonian dynamics\\ on interacting
quantum systems}

\vfill Pawel Wocjan, Dominik Janzing, Martin R{\"o}tteler, Thomas Beth

{\small Institut f"ur Algorithmen und Kognitive Systeme}

{\small Universit"at Karlsruhe}
\end{center}
\vfill
\end{slide}

\begin{slide}
{\large {\bf Control-theoretical model of quantum computers}}
\vfill
{\blue Time evolution} defined by the {\blue
time-dependent Schr\"odinger equation}
\[
\frac{d}{dt} U(t) = -i H(t) U(t)\,,\quad U(0)=\onemat
\]
$H(t)$ Hamiltonian 

$U(t)$ is the resulting unitary after time $t$
\vfill
\[
H(t) = H_d + \sum_{i=1}^m a_i(t) H_i
\]
$H_d$ {\blue drift Hamiltonian} internal to the system

$\sum_{i=1}^m a_i(t) H_i$ {\blue control Hamiltonian} externally
changeable
\vfill
{\red Complexity of a unitary $U$}: minimal time required to implement
$U$
\vfill
\end{slide}

\begin{slide}
{\large {\bf Fast control limit}}
\vfill
strength of the control Hamiltonians arbitrarily large $\Longrightarrow$

{\blue control unitaries in control group $K$} arbitrarily fast
\vfill
{\red Complexity of a unitary $U$}: minimal time
\[
U=\exp(-i H t_N) V_N \cdots \exp(-i H t_2) V_2 \exp(-i H t_1) V_1
\]
\vfill
express it as the evolution according to the piecewise constant Hamiltonian
\[
U = U_N \exp(-i U_N^\dagger H U_N t_N) \cdots \exp(-i U_2^\dagger H U_2 t_2)
    \exp(-i U_1^\dagger H U_1 t_1)
\]
\vfill
\end{slide}


\begin{slide}
{\large {\bf Average Hamiltonian Theory}}
\vfill
{\blue Dyson expansion}
\[
U(t)={\cal T} \exp(\int_{\tau=0}^t d\tau H(\tau))=\exp(-i\bar{H}t)
\]
\vfill
{\blue Magnus expansion} $\bar{H}=\bar{H}^{(0)}+\bar{H}^{(1)}+\cdots$
\begin{eqnarray*}
{\red\bar{H}^{(0)}} & {\red =} & {\red \frac{1}{t}\int_0^t d\tau
H(\tau)\,,\quad\quad\quad\quad\ \quad\quad\quad\quad
\|\bar{H}^{(0)}\|\le 1} \\
\bar{H}^{(1)} & = & \frac{-i}{2t}\int_0^t
d\tau' \int_0^t\tau'' [H(\tau'),H(\tau'')]\,,\quad
\|\bar{H}^{(1)}\|\le t/2
\end{eqnarray*}
\vfill
{\red $\Longrightarrow$ $U(t)$ essentially determined by $\bar{H}^{(0)}$ for
small $t$}
\vfill
\end{slide}


\begin{slide}
{\large {\bf Simulating Hamiltonians (first order simulation)}} 

\vfill
{\blue $H$ can simulate $\tilde{H}$ with overhead $1$}, written
$\tilde{H}\le H$, and {\blue $N$ time steps} iff
\[
\tilde{H}=\sum_{j=1}^N p_j U_j^\dagger H U_j\,,
\]
where $p_1,\ldots,p_N$ are positive real numbers summing to $1$ and
$U_j\in K$ are control operations, i.e.
{\red $\tilde{H}$ is a convex sum of unitary conjugates of $H$}
\vfill
$H$ can simulate $\tilde{H}$ with overhead {\red $\tau$} iff 
$\tilde{H} \le {\red \tau} H$ 
\vfill 
{\red time-optimal simulation} problem $\longrightarrow$
{\red convex optimization} problem 
\vfill
\end{slide}

\begin{slide}
{\large {\bf Composite systems}}
\vfill
system consisting of {\blue $n$ qudits}  
$\quad {\cal H}=\C^d\otimes\C^d\otimes\cdots\otimes\C^d$ 
\vfill
{\blue pair-interaction Hamiltonian}
\[
H = 
\sum_{k<l} \sum_{\alpha\beta} J_{kl;\alpha\beta} 
\sigma_\alpha^{(k)}\sigma_\beta^{(l)} +
\sum_k \sum_{\alpha} r_{k;\alpha} \sigma_\alpha^{(k)}\,,
\]
where $B:=\{\sigma_\alpha\mid\alpha=1,\ldots,d^2-1\}$ is a basis of $su(d)$

%{\blue coupling matrix}
%\[
%J=\left(
%\begin{array}{c|c|c|c|c}
%0      & J_{12} & J_{13} & \cdots & J_{1n} \\ \hline
%J_{21} & 0      & J_{23} & \cdots & J_{2n} \\ \hline
%\vdots & \vdots &        &        &        \\ \hline
%J_{n1} & J_{n2} & J_{n3} &        & 0 
%\end{array}
%\right)\in\R^{(d^2-1)\,n\times (d^2-1)\,n}
%\]

\vfill
{\blue control operations} are all {\red local} operations in
\[
K:=SU(d)\otimes SU(d)\otimes\cdots\otimes SU(d)
\]
\vfill
\end{slide}

\begin{slide}
{\blue Decoupling/Selective decoupling}

\quad quantum computing requires the ability to perform gates

\quad couplings created by gates can only originate from natural couplings

\quad naturally available Hamiltonians do not couple specific pairs of qubits

\quad {\red $\longrightarrow$ turn off disturbing couplings}

\vfill
{\blue Time-reversal}

\quad evolution of a single spin seems is phenomenologically a relaxation/

\quad dephasing process

%\quad (described by Bloch equations)

\quad however, it is not always irreversible

\quad {\red $\longrightarrow$ refocusing}\quad simulate $-H$ when $H$ is
the system Hamiltonian

\vfill 
{\blue Universal simulation of quantum dynamics}

\quad simulating quantum dynamics on a classical computer is hard

\quad {\red $\longrightarrow$ universal quantum simulator}

\quad\quad\quad use a quantum computer to simulate arbitrary dynamics
\vfill
\end{slide}

\begin{slide}
{\bf Mathematical tools}

{\blue Vector majorization}

let $x=(x_1,x_2,\ldots,x_d)\in\R^d$

define
$x^\downarrow=(x^\downarrow_1,x^\downarrow_2,\ldots,x^\downarrow_d)\in\R^d$
as the vector whose components are the same as of $x$, but ordered so that
\[
x^\downarrow_1\ge x^\downarrow_2\ge\ldots\ge x^\downarrow_d
\]

{\blue $x$ majorizes $y$}, written {\blue $x\preceq y$}, if
\[
\sum_{j=1}^k x^\downarrow_j \le \sum_{j=1}^k y^\downarrow_j
\]
for $k=1,\ldots,d-1$ and 
\[
\sum_{j=1}^d x^\downarrow_j = \sum_{j=1}^d y^\downarrow_j
\]
\end{slide}

\begin{slide}
{\bf Mathematical tools} 

{\blue Operator majorization}
\vfill
let $A$ and $B$ be two Hermitian matrices of size $d$

we say {\blue $A$ majorizes $B$}, written {\blue $A\preceq B$}, if
\[
\lambda(A)\preceq\lambda(B)\,,
\]
where $\lambda(A)$ and $\lambda(B)$ are the spectra (the vectors of
eigenvalues) of $A$ and $B$, respectively.
\vfill
{\bf Uhlmann's theorem} 

we have $A\preceq B$ iff
\[
A=\sum_{j=1}^N p_j U_j^\dagger B U_j\,,
\]
where $\{p_j\}$ is a probability distribution and $\{U_j\}$ is a set
of unitaries
\vfill
\end{slide}

\end{document}









